\chapter[Introdução]{Introdução}
%\addcontentsline{toc}{chapter}{Introdução}
% ----------------------------------------------------------
O tema central deste Trabalho de Conclusão de Curso (TCC) é a aplicação de Inteligência Artificial Generativa na automação do processo de correção de avaliações acadêmicas. A delimitação do tema foca especificamente no desenvolvimento de um artefato tecnológico (aplicativo móvel e sistema web) voltado para a correção de provas objetivas — aquelas compostas exclusivamente por questões de múltipla escolha com alternativas fechadas.

Esta delimitação é estratégica e necessária para garantir a viabilidade técnica e a precisão do artefato no estágio atual da tecnologia. Embora a IA Generativa já apresente avanços na correção de questões dissertativas, a natureza subjetiva dessas avaliações exige padrões predefinidos complexos e maior propensão a alucinações por parte do modelo, o que poderia comprometer a confiabilidade inicial do sistema. Ao focar em questões objetivas, o projeto busca atingir um nível de precisão próximo ao "erro zero" na extração de dados, permitindo que a inovação se concentre na flexibilidade da entrada (fotos de qualquer formato de prova) e na usabilidade para o docente.

\section{Tema e Delimitação}\label{justif}

Tema e Delimitação...

\section{Problema}\label{justif}

Os trabalhos desenvolvidos...

\section{Justificativa}\label{objgeral}

Nas seções abaixo estão descritos o objetivo geral e os objetivos específicos deste TCC.

\subsection{Objetivo geral}\label{objger}

O objetivo geral deste trabalho é...

\subsection{Objetivos específicos}\label{objesp}

Os objetivos específicos são...

\section{Organização do trabalho}\label{orgtra}

A introdução termina com a apresentação dos demais capítulos do trabalho. O capítulo 2 contém o referencial teórico deste trabalho. No capítulo 3 é apresentado o desenvolvimento, o cenário, os testes e a discussão dos resultados. Por fim, temos a conclusão, onde discutiremos os resultados, as dificuldades encontradas e faremos sugestões de trabalhos futuros.

Este trabalho segue as normas estabelecidas pela ABNT \cite{NBR14724:2011}, conforme orientações do modelo canônico \cite{abntex2modelo}.